\begin{figure}[t]
\centering
\begin{tikzpicture}
\begin{axis}[txaiaxis,
  xmode=log, log basis x=10,
  xmin=0.85, xmax=2400, ymin=-0.03, ymax=1.06,
  ytick={0,0.25,0.5,0.75,1},
  xlabel={Features deleted in attribution order, $k$},
  ylabel={Mean detector confidence},
  legend style={at={(0.02,0.02)},anchor=south west}]
\addplot[black!45,dashed,forget plot,line width=0.5pt,mark=none,domain=0.85:2400]
  {0.994894};
\addplot[black!45,dotted,forget plot,line width=0.5pt,mark=none,domain=0.85:2400]
  {0.5};
\node[anchor=east,font=\tiny,text=black!60,inner sep=1.5pt]
  at (axis cs:2300,0.93) {undeleted mean};
\node[anchor=west,font=\tiny,text=black!60,inner sep=1.5pt]
  at (axis cs:1.05,0.55) {decision boundary};
\addplot[color=oiblue,mark=*] coordinates {(1,0.972881) (2,0.877338) (5,0.572955) (10,0.413362) (20,0.220953) (50,0.0340235) (100,0.00649059) (200,0.00182138) (500,0.000478727) (1000,0.00024928) (2000,0.00024928)};
\addlegendentry{TreeSHAP}
\addplot[color=oiorange,mark=square*] coordinates {(1,0.980855) (2,0.932603) (5,0.807352) (10,0.741458) (20,0.679874) (50,0.621245) (100,0.562849) (200,0.485951) (500,0.355795) (1000,0.155096) (2000,0.00413824)};
\addlegendentry{LIME}
\addplot[color=black,mark=diamond*] coordinates {(1,0.994894) (2,0.994894) (5,0.994761) (10,0.994782) (20,0.994685) (50,0.978326) (100,0.969866) (200,0.968095) (500,0.487827) (1000,0.338584) (2000,0.00361998)};
\addlegendentry{constant}
\addplot[color=oipurple,mark=triangle*] coordinates {(1,0.994866) (2,0.994862) (5,0.994761) (10,0.994498) (20,0.994198) (50,0.992154) (100,0.987783) (200,0.971832) (500,0.864079) (1000,0.528456) (2000,0.0351706)};
\addlegendentry{random}
\end{axis}
\end{tikzpicture}
\caption{Deletion behaviour of the four explanation maps on the $200$ alerts of
the single-corpus instantiation of Section~\ref{sec:numsetup}: mean confidence
of the \textsf{histgb} detector after deleting the $k$ highest-attributed
features. The two grey guides mark the undeleted mean confidence
($0.9949$) and the decision boundary. Read the
curves for their \emph{ordering and shape}, not for a per-alert guarantee: only
TreeSHAP drives the detector below the boundary within the first ten features,
while the constant control tracks the undeleted level through $k=200$ because its
fixed ordering is uninformative rather than adversarial. LIME's curve is
measured on the reduced grid its cost forced (Section~\ref{sec:numlimits}) and
is not directly comparable to the others. Single corpus, single detector, single
seed; the ordering is not asserted beyond this instantiation.}
\label{fig:deletion}
\end{figure}

\begin{figure}[t]
\centering
\begin{tikzpicture}
\begin{groupplot}[txaiaxis,
  group style={group size=1 by 2, vertical sep=5pt,
                x descriptions at=edge bottom},
  % ``scale only axis'' so that the two heights below size the plotting boxes
  % themselves: without it pgfplots subtracts the shared x labels from the short
  % lower panel and aborts with a negative plot height. The width then excludes
  % the y label and tick labels, hence 0.82 rather than 0.97:
  % at 0.84 the picture runs 2.4pt past the column.
  scale only axis, width=0.82\columnwidth,
  xmode=log, log basis x=10,
  xtick={0.001,0.005,0.01,0.05},
  xticklabels={$0.1\%$,$0.5\%$,$1\%$,$5\%$},
  x tick label style={/pgf/number format/assume math mode=true},
  xmin=0.0008, xmax=0.065]
\nextgroupplot[height=0.42\columnwidth, ymin=0.835, ymax=1.025,
  ytick={0.85,0.9,0.95,1},
  ylabel={Robustness score $B_{\Gamma,r}$},
  legend style={at={(0.03,0.05)},anchor=south west}]
\addplot[color=oiblue,mark=*] coordinates {(0.001,0.979518) (0.005,0.959232) (0.01,0.946481) (0.05,0.900699)};
\addlegendentry{\textsf{append\_bytes}}
\addplot[color=oigreen,mark=square*] coordinates {(0.001,0.990068) (0.005,0.971245) (0.01,0.955829) (0.05,0.906005)};
\addlegendentry{\textsf{add\_imports}}
\addplot[color=oisky,mark=triangle*] coordinates {(0.001,0.984141) (0.005,0.957147) (0.01,0.936349) (0.05,0.854312)};
\addlegendentry{\textsf{generic}}
\addplot[color=black,mark=diamond*,dashed] coordinates {(0.001,1) (0.005,1) (0.01,1) (0.05,1)};
\addlegendentry{constant control}
\nextgroupplot[height=0.09\columnwidth, ymin=0.572, ymax=0.5894,
  ytick={0.575}, yticklabels={$0.575$},
  xlabel={Perturbation budget $r$ (fraction of the input)}]
\addplot[color=oipurple,mark=triangle*,dashed] coordinates {(0.001,0.574958) (0.005,0.575177) (0.01,0.575433) (0.05,0.575207)};
\node[anchor=north west,font=\scriptsize,inner sep=1.5pt]
  at (rel axis cs:0.02,0.98) {random control};
\end{groupplot}
\draw[black,line width=0.4pt] ([shift={(-2.2pt,-1.6pt)}]group c1r1.south west) -- ([shift={(2.2pt,1.6pt)}]group c1r1.south west);
\draw[black,line width=0.4pt] ([shift={(-2.2pt,-1.6pt)}]group c1r1.south east) -- ([shift={(2.2pt,1.6pt)}]group c1r1.south east);
\draw[black,line width=0.4pt] ([shift={(-2.2pt,-1.6pt)}]group c1r2.north west) -- ([shift={(2.2pt,1.6pt)}]group c1r2.north west);
\draw[black,line width=0.4pt] ([shift={(-2.2pt,-1.6pt)}]group c1r2.north east) -- ([shift={(2.2pt,1.6pt)}]group c1r2.north east);
\end{tikzpicture}
\caption{Robustness score \eqref{eq:robustness} against the adversary's budget,
each point a mean over $500$ alerts $\times$ $20$ sampled transformations. The
three solid curves are TreeSHAP under the three admitted transformation
families; the two dashed lines are the controls, both measured under
\textsf{append\_bytes}. The vertical axis is broken at the marked corners
because the random control lies far below everything else, and on one continuous
scale it flattens the three curves that actually move into the top seventh of
the panel. Section~\ref{sec:numrobust} reads the three movements and says what
the flat lines do and do not license. LIME is absent: its cost admitted only the
single $1\%$ point, reported in the text. Single corpus, single detector, single
seed.}
\label{fig:robustness}
\end{figure}

\begin{figure}[t]
\centering
\begin{tikzpicture}
\begin{axis}[txaiaxis,
  xmin=-0.03, xmax=1.03, ymin=-0.05, ymax=1.16,
  ytick={0,0.25,0.5,0.75,1},
  xlabel={Disclosure threshold $\tau_d$ (at $\tau_f=0.5$, $\tau_b=0.8$)},
  ylabel={Fraction of alerts admissible},
  legend style={at={(0.03,0.70)},anchor=west}]
\addplot[color=black,mark=diamond*] coordinates {(0,0) (0.1,0) (0.2,0) (0.3,0) (0.4,0) (0.5,0) (0.6,0) (0.7,1) (0.8,1) (0.9,1) (1,1)};
\addlegendentry{constant}
\addplot[color=oiblue,mark=*] coordinates {(0,0) (0.1,0) (0.2,0) (0.3,0) (0.4,0) (0.5,0) (0.6,0.136) (0.7,0.3) (0.8,0.944) (0.9,1) (1,1)};
\addlegendentry{TreeSHAP}
\addplot[color=oiorange,mark=square*] coordinates {(0,0) (0.1,0) (0.2,0) (0.3,0.01) (0.4,0.03) (0.5,0.03) (0.6,0.03) (0.7,0.03) (0.8,0.03) (0.9,0.03) (1,0.03)};
\addlegendentry{LIME}
\addplot[color=oipurple,mark=triangle*] coordinates {(0,0) (0.1,0) (0.2,0) (0.3,0) (0.4,0) (0.5,0) (0.6,0) (0.7,0) (0.8,0) (0.9,0) (1,0)};
\addlegendentry{random}
\draw[<->,black!55,line width=0.5pt]
  (axis cs:0.7,0.34) --
  (axis cs:0.7,0.96);
\node[anchor=south,font=\scriptsize,text=black!70,inner sep=2pt]
  at (axis cs:0.7,1.02)
  {$\tau_d=0.7$: $1.00$ vs $0.30$};
\end{axis}
\end{tikzpicture}
\caption{The four numeric conditions of Definition~\ref{def:admissible}, applied
to $200$ alerts as the disclosure threshold is relaxed with the other two
thresholds held at the values used throughout Section~\ref{sec:numerical} (the
analyst role; the other three roles coincide with it on this slice). This is the
measurement behind the calibration argument, and the marked span is the whole of
it: the constant control becomes admissible on \emph{every} alert at a
disclosure threshold where TreeSHAP is still admissible on under a
third, so on this slice the inequalities alone accept a vector that
explains nothing over a strictly wider region than they accept a Shapley
attribution. The gap is a property of this instantiation, not a theorem; what
generalises is that the numeric conditions cannot by themselves exclude the
degenerate map, which is what \eqref{eq:calibration} is for. Single corpus,
single detector, single seed.}
\label{fig:admissible}
\end{figure}

\begin{figure}[t]
\centering
\begin{tikzpicture}
\begin{axis}[txaiaxis,
  xmode=log, ymode=log, log basis x=10, log basis y=10,
  xmin=1.2e-3, xmax=0.9, ymin=2.5e-5, ymax=0.9,
  xlabel={Right-hand side of \eqref{eq:bridge}},
  ylabel={Measured $\mathrm{Adv}^{\mathrm{TV}}_{E,\Gamma}$},
  legend style={at={(0.03,0.88)},anchor=west}]
\addplot[black,densely dashed,mark=none,domain=1.2e-3:0.9,samples=2,
  forget plot] {x};
\addplot[only marks,color=oipurple,mark=o,mark size=1.3pt] coordinates
  {(0.0204822,0.00318977) (0.0204822,0.00147003) (0.0204822,0.000428206) (0.0407676,0.00808549) (0.0407676,0.0037701) (0.0407676,0.00108854) (0.0535191,0.0117031) (0.0535191,0.00545006) (0.0535191,0.00157379) (0.0993008,0.0240962) (0.0993008,0.0117472) (0.0993008,0.00334018) (0.00993218,0.000697665) (0.00993218,0.000310124) (0.00993218,8.93398e-05) (0.0287551,0.00227224) (0.0287551,0.000846422) (0.0287551,0.000242376) (0.0441712,0.00433315) (0.0441712,0.00152684) (0.0441712,0.000411268) (0.0939951,0.01637) (0.0939951,0.00577858) (0.0939951,0.0015979) (0.0158594,0.000283066) (0.0158594,0.000146253) (0.0158594,4.77275e-05) (0.0428526,0.000680479) (0.0428526,0.000447587) (0.0428526,0.000150986) (0.0636506,0.00105588) (0.0636506,0.000628775) (0.0636506,0.000220746) (0.145688,0.00637432) (0.145688,0.00352902) (0.145688,0.00115206) (0.404867,0.00264179) (0.404867,0.00113381) (0.404867,0.00030947) (0.425042,0.00105139) (0.425042,0.000347172) (0.425042,8.86905e-05) (0.424823,0.0011673) (0.424823,0.000482225) (0.424823,0.0001249) (0.424567,0.00100996) (0.424567,0.000391302) (0.424567,0.000106794) (0.424793,0.00115771) (0.424793,0.000342361) (0.424793,9.15705e-05) (0.424545,0.00097378) (0.424545,0.000327403) (0.424545,8.7806e-05) (0.42483,0.00110201) (0.42483,0.000450366) (0.42483,0.000120711) (0.424868,0.00127503) (0.424868,0.000371161) (0.424868,0.000114166) (0.424822,0.00108001) (0.424822,0.000411565) (0.424822,0.000111149) (0.424613,0.000943207) (0.424613,0.000294661) (0.424613,8.67865e-05) (0.424929,0.000902239) (0.424929,0.000385552) (0.424929,0.000105812) (0.424849,0.000839363) (0.424849,0.000312927) (0.424849,7.85972e-05) (0.424834,0.00142645) (0.424834,0.000497975) (0.424834,0.000132481)};
\addlegendentry{generic constant $L_\pi=1$}
\addplot[only marks,color=oiblue,mark=x,mark size=1.9pt] coordinates
  {(0.0201191,0.00318977) (0.012402,0.00147003) (0.00394827,0.000428206) (0.0400448,0.00808549) (0.0246848,0.0037701) (0.00785859,0.00108854) (0.0525702,0.0117031) (0.0324059,0.00545006) (0.0103166,0.00157379) (0.0975402,0.0240962) (0.0601267,0.0117472) (0.0191418,0.00334018) (0.00975608,0.000697665) (0.00601394,0.000310124) (0.00191458,8.93398e-05) (0.0282453,0.00227224) (0.0174112,0.000846422) (0.005543,0.000242376) (0.043388,0.00433315) (0.0267457,0.00152684) (0.00851469,0.000411268) (0.0923285,0.01637) (0.0569141,0.00577858) (0.018119,0.0015979) (0.0155782,0.000283066) (0.00960289,0.000146253) (0.00305715,4.77275e-05) (0.0420928,0.000680479) (0.0259473,0.000447587) (0.00826051,0.000150986) (0.062522,0.00105588) (0.0385405,0.000628775) (0.0122696,0.000220746) (0.143105,0.00637432) (0.0882142,0.00352902) (0.0280836,0.00115206) (0.397689,0.00264179) (0.245148,0.00113381) (0.0780445,0.00030947) (0.417505,0.00105139) (0.257363,0.000347172) (0.0819334,8.86905e-05) (0.417291,0.0011673) (0.257231,0.000482225) (0.0818913,0.0001249) (0.417039,0.00100996) (0.257076,0.000391302) (0.0818419,0.000106794) (0.417262,0.00115771) (0.257213,0.000342361) (0.0818856,9.15705e-05) (0.417018,0.00097378) (0.257062,0.000327403) (0.0818377,8.7806e-05) (0.417298,0.00110201) (0.257235,0.000450366) (0.0818927,0.000120711) (0.417335,0.00127503) (0.257258,0.000371161) (0.0819,0.000114166) (0.41729,0.00108001) (0.25723,0.000411565) (0.0818912,0.000111149) (0.417084,0.000943207) (0.257103,0.000294661) (0.0818508,8.67865e-05) (0.417395,0.000902239) (0.257295,0.000385552) (0.0819116,0.000105812) (0.417316,0.000839363) (0.257246,0.000312927) (0.0818962,7.85972e-05) (0.417301,0.00142645) (0.257237,0.000497975) (0.0818933,0.000132481)};
\addlegendentry{Dobrushin coefficient $\delta$}
\end{axis}
\end{tikzpicture}
\caption{Theorem~\ref{thm:bridge} on the instantiated kernel. Each of the
$75$ plotted configurations lies below the diagonal $y=x$ (dashed), so the
inequality is not violated, and each lies far below it, so the bound is
loose. Replacing the generic constant $L_\pi=1$ by the ergodic
coefficient \eqref{eq:dobrushin} moves each point left without moving it above
the diagonal, which is the sharpening claimed in the text. The remaining $36$ of the $111$ recomputed
configurations are the constant control's, omitted because both sides are
exactly zero and so have no position on logarithmic axes; they are also the only
configurations in which the bound is attained. Points are configurations, not alerts, and
holding on the configurations of one instantiation is not a claim that the
bound is tight or loose in general.}
\label{fig:bridge}
\end{figure}

\begin{table}[t]
\caption{Strong Stackelberg equilibrium of the $15\times13$ release game
of Section~\ref{sec:numgame}, over the two weights the utilities of
Definition~\ref{def:game} leave free: $\eta$ on evidence failure and $\rho$ on
disclosure. $k^\star$ is the equilibrium release width on the lattice
$\{5,10,20,50,100\}$ and $U_D$ the defender's equilibrium value. The width is
nondecreasing in $\eta$ and nonincreasing in $\rho$: an evidence obligation is
the only term that pays for release, and at $\eta=0$ the equilibrium releases
the least the lattice allows at every $\rho$. The equilibrium is pure at all
$18$ settings---the gain from randomizing does not
exceed $4\times10^{-16}$---and the attacker plays pure observation
($M_4$) at every one of them.}
\label{tab:numgame}
\centering
\footnotesize
\begin{tabular}{@{}lcccccc@{}}
\toprule
& \multicolumn{2}{c}{$\eta=0$} & \multicolumn{2}{c}{$\eta=1$} & \multicolumn{2}{c}{$\eta=4$}\\
$\rho$ & $k^\star$ & $U_D$ & $k^\star$ & $U_D$ & $k^\star$ & $U_D$\\
\midrule
$0$ & $5$ & $-0.284$ & $100$ & $-0.434$ & $100$ & $-0.557$\\
$0.25$ & $5$ & $-0.313$ & $100$ & $-0.617$ & $100$ & $-0.740$\\
$0.5$ & $5$ & $-0.342$ & $100$ & $-0.799$ & $100$ & $-0.922$\\
$1$ & $5$ & $-0.400$ & $50$ & $-1.138$ & $100$ & $-1.287$\\
$2$ & $5$ & $-0.516$ & $5$ & $-1.377$ & $100$ & $-2.017$\\
$4$ & $5$ & $-0.748$ & $5$ & $-1.609$ & $50$ & $-3.415$\\
\bottomrule
\end{tabular}
\end{table}
